\PassOptionsToPackage{dvipsnames,table}{xcolor}
\documentclass[11pt,a4paper]{article}

\usepackage{Act}
\begin{document}
\input{\detokenize{/home/fenarius/Travail/Cours/cpge-info/latex/Macros.tex}}
\ModeExercice
\setboolean{corrige}{false}
\ICX{Devoir ITC}{2}{pcsi}{Bases de Python}{10}
\setcounter{Exercise}{0}


\begin{Exercise} \points{10}\\
    Une réponse brève d'une ligne est attendue dans le cadre qui suit immédiatement la question, on ne \textit{demande pas} de justification.
	\Question{En Python, quelle est la valeur de la variable {\tt a} après exécution de l'instruction \mintinline{python}{a = (16 // 5)**3 - 2}}
	\rep{0}
	\tcor{\tt 25}

	\Question{Si la variable {\tt x} vaut 42, quelle est la valeur de \mintinline{python}{x % 7 == 0 and (x < 30 or x % 2 == 0)} ?}
	\rep{0}
	\tcor{\tt True}

	\Question{Si {\tt lst} vaut {\tt [3, 7, 12, 18, 25]}, que vaut \mintinline{python}{lst[-1] - lst[1] + len(lst)} ?}
	\rep{0}
	\tcor{\tt 23}

	\Question{Pour quelle(s) valeurs de la variable {\tt i} sera effectuée la boucle \mintinline{python}{for i in range(20, 5, -4)}}
	\rep{0}
	\tcor{{\tt 20, 16, 12, 8}}

	\Question{Si {\tt s} est la chaine de caractères {\tt "Python"} quel est l'affichage produit par \mintinline{python}{print(s[:2] * 2 + s[-1])}}
	\rep{0}
	\tcor{{\tt "PyPyn"}}

	\Question{Écrire en une seule instruction Python l'affectation permettant d'échanger les valeurs de deux variables {\tt a} et {\tt b} sans variable intermédiaire.}
	\rep{0}
	\tcor{\mintinline{python}{a, b = b, a}}

	\Question{Expliquer l'origine de l'erreur {\tt TypeError: 'tuple' object does not support item assignment} en Python.}
	\rep{0}
	\tcor{Cela signifie qu'on a tenté de modifier un élément d'un tuple, or le type {\tt tuple} n'est pas mutable (il est immuable).}

	\Question{Écrire l'instruction conditionnelle permettant de tester si un entier {\tt n} est compris entre 10 et 50 (inclus) et est divisible par 3.}
	\rep{0}
	\tcor{{\tt if 10 <= n <= 50 and n \% 3 == 0:}}

	\Question{Écrire une instruction permettant de créer \textit{par compréhension} la liste des carrés des entiers de 1 à 10 inclus : {\tt [1, 4, 9, 16, 25, 36, 49, 64, 81, 100]}}
	\rep{0}
	\tcor{\mintinline{python}{l = [k**2 for k in range(1, 11)]}}

	\Question{Quel sera le contenu de la liste {\tt a} après exécution des instructions suivantes :
		\inputpartPython{qcm2.py}{}{}{1}{4}
	}
	\rep{0}
	\tcor{\mintinline{python}{[10, 99, 30, 40]}}
\end{Exercise}

\newpage
\begin{Exercise}[title={QCM}] \points{10}\\
	Dans cet exercice, une question peut avoir \textit{zéro, une ou plusieurs bonnes réponses}. Pour chaque question, cocher les cases correspondantes aux bonnes réponses.
	\Question{Que peut-on dire du résultat de l'expression \mintinline{python}{15 // 4 + 7 / 2} ?}\\
	\begin{tabularx}{\linewidth}{@{}XXXX@{}}
		\bc C'est un {\tt int} & \gc C'est un {\tt float} & \bc Sa valeur est 6 & \gc Sa valeur est 6.5 \\
	\end{tabularx}
    \Question{Quelles sont les propositions exactes concernant les fonctions en Python ?}
    \begin{itemize}
    \item[\gc] Une fonction sans instruction {\tt return} explicite renvoie la valeur {\tt None}
    \item[\bc] Les arguments d'une fonction doivent obligatoirement être du même type
    \item[\gc] Une fonction peut renvoyer plusieurs valeurs séparées par des virgules
    \item[\gc] Les variables définies à l'intérieur d'une fonction ont une portée locale à cette fonction
    \end{itemize}
	\Question{Quel(s) test(s) sont vrais si et seulement si l'entier {\tt n} est impair ?}\\
	\begin{tabularx}{\linewidth}{@{}XXXXX@{}}
		\gc {\tt n\%2!=0} & \gc {\tt n\%2==1} & \bc {\tt n//2==1} & \bc {\tt 2\%n==1} & \gc {\tt not(n\%2==0)} \\
	\end{tabularx}
	\Question{Si {\tt lst} est une liste (type {\tt list}), cocher les instructions valides (qui ne provoquent pas d'erreur)} \\
	\begin{tabularx}{\linewidth}{@{}XXXXX@{}}
		\gc {\tt lst + [3]} & \bc {\tt lst + 3} & \gc {\tt lst * 2} & \gc {\tt lst[0] = 1} & \bc {\tt lst(0)} \\
	\end{tabularx}
	\Question{Parmi les types de Python suivants, lesquels sont \textit{mutables} (modifiables en place) ?}\\
	\begin{tabularx}{\linewidth}{@{}XXXXX@{}}
		\bc {\tt int} & \bc {\tt float} & \bc {\tt str} & \gc {\tt list} & \bc {\tt tuple} \\
	\end{tabularx}
	\Question{On suppose que {\tt x} est un entier valant 8, quelles expressions seront évaluées à {\tt True} ?}\\
	\begin{tabularx}{\linewidth}{@{}XXXX@{}}
		\gc {\tt x>5 and x<10} & \gc {\tt not(x!=8)} & \bc {\tt x\%4==0 and x\%3==0} & \gc {\tt 10>x>=8} \\
	\end{tabularx}
	\Question{On suppose que {\tt s = 0} a été exécuté. Parmi les boucles suivantes, lesquelles permettent de calculer dans {\tt s} la somme des éléments de {\tt lst} ?} \\
	\begin{minipage}[b][3 cm][t]{7 cm}
        \begin{itemize}
		\item[\gc \ \ ]\ipp{qcm2.py}{8}{9}
		\item[\bc \ \ ]\ipp{qcm2.py}{11}{12}
		\end{itemize}
	\end{minipage}\hfill
	\begin{minipage}[b][3 cm][t]{7 cm}
		\begin{itemize}
		\item[\gc \ \ ]\ipp{qcm2.py}{14}{15}
		\item[\bc \ \ ]\ipp{qcm2.py}{17}{18}
		\end{itemize}
	\end{minipage}\par
    \Question{Quelles sont les affirmations vraies concernant le type {\tt str} de Python ?}
    \begin{itemize}
    \item[\bc] On peut modifier directement un caractère d'une chaine par affectation : {\tt s[0] = 'a'}
    \item[\gc] On peut concaténer deux chaines de caractères à l'aide de l'opérateur {\tt +}
    \item[\gc] La fonction {\tt len(s)} renvoie le nombre de caractères de la chaine {\tt s}
    \item[\bc] Une chaine de caractères ne peut pas être vide
    \item[\gc] L'instruction {\tt s = ""} permet de définir une chaine vide
    \end{itemize}
    \Question{Si {\tt lst} est la liste {\tt [10, 20, 30, 40, 50]}, quelles tranches permettent d'obtenir la liste {\tt [20, 30]} ?} \\
    \begin{tabularx}{\linewidth}{@{}XXXX@{}}
		\gc {\tt lst[1:3]} & \bc {\tt lst[1:2]} & \gc {\tt lst[-4:-2]}  & \gc {\tt lst[1:-2]} \\
	\end{tabularx}
    \Question{Après exécution de l'instruction {\tt l = [3*i for i in range(1, 6)]}, quelles affirmations concernant {\tt l} sont vraies ?}
    \begin{itemize}
    \item[\gc] {\tt len(l)} vaut 5 
    \item[\bc] {\tt 0} est l'un des éléments de {\tt l} 
    \item[\gc] {\tt l} est la liste {\tt [3, 6, 9, 12, 15]} 
    \item[\gc] Le dernier élément de {\tt l} est 15
    \end{itemize}
\end{Exercise}

\newpage
\begin{Exercise}[title = {Définir une fonction}]\\
    Écrire en Python une fonction {\tt nb\_pairs} qui prend en argument une liste d'entiers {\tt lst} et renvoie le nombre d'éléments pairs contenus dans cette liste. Par exemple, {\tt nb\_pairs([2, 5, 8, 11, 14])} doit renvoyer 3, et {\tt nb\_pairs([1, 3, 7])} doit renvoyer 0.
    \reponse{9}{5}
	\ifcorrige
	\corpartPython{qcm2.py}{}{}{20}{25}
	\fi
\end{Exercise}


\begin{Exercise}[title = {Exercice bonus}]\\
	Écrire une fonction {\tt deuxmax} qui prend en argument une liste d'entiers contenant au moins deux éléments et qui renvoie les deux plus grands éléments de cette liste. Par exemple {\tt deuxmax([14, -2, 6, 20, 8, 3])} renvoie {\tt 14, 20}.
    \rep{14}
	\ifcorrige
	\corpartPython{qcm2.py}{}{}{27}{41}
	\fi
\end{Exercise}

\end{document}
